\documentclass[fleqn,reqno,12pt]{article}

\usepackage{CSPArticleStyle}

% - - - - - - - - - - - - - - - - - - - - - - - -
% additional packages
% - - - - - - - - - - - - - - - - - - - - - - - -

\usepackage[margin=2.5cm]{geometry}


\bibliography{~/Desktop/icloud/MyRefGlobal.bib}


\title{The cognitive turn in computational modeling for experimental pragmatics}
\author{Michael Franke}
\date{}

\begin{document}
\maketitle

\begin{abstract}
  concrete.
\end{abstract}


\section{Experimental pragmatics and computational modeling}

The linguistic subfield of pragmatics studies how language is used for efficient communication and how communicated meaning is modulated, if not created, in a particular context by systematic, context-general patterns of use.
Its origins lie in the Philosophy of Language, in particular the analytical tradition of philosophy, in work on topics like \emph{speech-acts} \citep{Austin1962:How-to-do-Thing,Searle1969:Speech-Acts:-An}, \emph{conversational implicature} \citep{Grice1975:Logic-and-Conve}, or \emph{presupposition} \citep{Stalnaker1973:Presupposition,Karttunen1973:Presuppositions}.
Much early work on these topics remained invested in a logical-formal approach \citep{Gazdar1979:Pragmatics:-Imp,Kadmon2001:Formal-Pragmati}, but the early 2000's brought a deep and empowering incision: a widely adopted and ultimately very successful \emph{experimental turn} \citep{NoveckSperber2004:Experimental-Pr,Noveck2018:Experimental-Pr}, in parallel to similar empirical re-orientations in other hitherto very analytically minded subfields of linguistics like formal syntax \citep{Sprouse2007:A-program-for-e,Sprouse2023:The-Oxf-Handb-Exp-Syntax} and formal semantics \citep{MeibauerSteinbach2011:Experimental-Pr,CumminsKatsos2019:The-Oxford-Hand}.

\begin{itemize}
  \item computational modeling
  \item computational pragmatics (unrelated to XPrag) existed -> explains stylized facts
  \item modeling for experimental pragmatics -> directly engage with the empirical data
\end{itemize}

\begin{itemize}
  \item modeling in general:
  \begin{itemize}
    \item all models are wrong; some are useful
    \item even if all models are wrong, we should prefer models that are less wrong in the sense of Isaak
    \item so how our (less) wrong models useful?
    \item while for data: garbage in — garbage out
    \item for models: more assumptions in — more conclusions out
    \begin{itemize}
      \item e.g., causal modeling
      \item reasonable assumptions in; interesting conclusions out
    \end{itemize}
  \end{itemize}
  \begin{itemize}
    \item all models contain some assumptions about the \textbf{data-generating process}
    \item more assumptions about the nature of that process → more information gained from the data
  \end{itemize}
  \item two specific contributions of explicit (computational / formal) modeling in XPrag:
  \begin{enumerate}
    \item \textbf{bridge to linguistic theory:} (a successful past)
    \begin{itemize}
      \item give linguistic theory (e.g., latent concepts) an explicit role in our picture of the data-generating process via theory-driven bespoke modeling
      \item experimental turn: experiments based on theory
      \item bring in the theory as a first-order citizen → theory-driven statistical modeling
    \end{itemize}
    \item \textbf{bridge to general cognition:} (a possible (bright) future)
    \begin{itemize}
      \item frame pragmatic reasoning in the broader picture of human cognition
      \item stuff for the future: proclaiming the *cognitive turn* in experimental pragmatics
    \end{itemize}
  \end{enumerate}
\end{itemize}

\newpage


\printbibliography

\end{document}
